\documentclass[11pt]{article}
\usepackage{amsmath, amssymb}

\title{Polynomials 2}
\date{}

\begin{document}
	
	\maketitle
	
	\begin{enumerate}
		
		\item \textbf{(O2) (Hensel's Lemma)}  
		Let $P \in \mathbb{Z}[x]$. Suppose that $p \mid P(x)$ and $p \nmid P'(x)$, where $P'$ is the derivative of $P$. Prove that:
		\begin{itemize}
			\item[(a)] There exists $y \in \mathbb{Z}$ such that $p^2 \mid P(y)$.
			\item[(b)] For each $n \in \mathbb{N}$, there exists $y \in \mathbb{Z}$ such that $p^n \mid P(y)$.
		\end{itemize}
		
		\item \textbf{(O2)}  
		Let $Q \in \mathbb{Z}[x]$ be a non-constant polynomial. Show that there exist arbitrarily large primes $p$ such that $p \mid Q(n)$ for some $n \in \mathbb{Z}$.
		
		\item \textbf{(O1)}  
		Let $P \in \mathbb{Z}[x]$ be irreducible. Show that the quantity
		\[
		\gcd\big(P(n), P'(n)\big)
		\]
		is bounded as $n$ ranges over $\mathbb{Z}$.
		
		\item \textbf{(O1)}  
		Let $P \in \mathbb{Z}[x]$.
		\begin{itemize}
			\item[(a)] Prove that if $P$ has a root in $\mathbb{Z}$, then $P$ has a root modulo $p$ for every prime $p$.
			\item[(b)] Prove that if there exists a prime $p$ such that $P$ has no root modulo $p$, then $P$ has no root in $\mathbb{Z}$.
		\end{itemize}
		
		\item \textbf{(O3)}  
		Let $P \in \mathbb{Z}[x]$ be non-constant. Does there necessarily exist a prime $p$ such that the sequence
		\[
		v_p(P(1)),\ v_p(P(2)),\ v_p(P(3)),\dots
		\]
		is unbounded?
		
	\end{enumerate}
	
\end{document}